\begin{table}[!htbp]
\setlength{\abovecaptionskip}{0pt}
\centering
\caption{Cross-Jurisdiction Heterogeneity in Administrative-Litigation Effects}
\label{tab:admin_cross_jurisdiction_heterogeneity}
\begin{threeparttable}
\begin{tabular}{lcccc}
\toprule
 & (1) & (2) & (3) & (4) \\
Subsample & Basic Court & Elevated Court & Local Plaintiff & Non-local Plaintiff \\
\midrule
Treatment $\times$ Post & 0.048$^{***}$ & -0.001 & 0.037$^{***}$ & 0.030$^{***}$ \\
& (0.005) & (0.008) & (0.005) & (0.007) \\
\addlinespace
Coefficient equality test ($p$) & \multicolumn{2}{c}{$<0.001$} & \multicolumn{2}{c}{0.344} \\
\addlinespace
Pre-treatment government win rate & 0.731 & 0.687 & 0.723 & 0.712 \\
Observations & 504,130 & 153,083 & 554,554 & 102,659 \\
$R^2$ & 0.070 & 0.056 & 0.070 & 0.086 \\
Plaintiff entity & Yes & Yes & Yes & Yes \\
Opposing counsel & Yes & Yes & Yes & Yes \\
Court Fixed Effects & Yes & Yes & Yes & Yes \\
Year Fixed Effects & Yes & Yes & Yes & Yes \\
Cause-Group Fixed Effects & Yes & Yes & Yes & Yes \\
\bottomrule
\end{tabular}
\begin{tablenotes}[flushleft]
\footnotesize
\item \textit{Note:} Linear-probability coefficients on Treatment $\times$ Post with the government-win indicator as the outcome, estimated on the indicated case sub-sample. Columns 1--2 split by court level: basic-level (district) people's courts versus intermediate, high, or specialized courts; the latter serve as a proxy for the cross-region adjudication arrangement of Liu, Wang, and Lyu (2023, \textit{Journal of Public Economics}). Columns 3--4 split by whether the plaintiff is local to the defendant city. The Coefficient equality test reports the two-sided $p$-value for $H_0$: column 1 coefficient = column 2 coefficient (and analogously for columns 3 vs 4) using the $z$-statistic computed from city- and court-clustered standard errors, treating the two sub-samples as independent and ignoring any residual within-city dependence across the split samples. Standard errors clustered by city and court. $^{*}p<0.10$, $^{**}p<0.05$, $^{***}p<0.01$.
\end{tablenotes}
\end{threeparttable}
\end{table}
